\documentclass[11pt,a4paper]{article}

\usepackage[utf8]{inputenc}
\usepackage[T1]{fontenc}
\usepackage{geometry}
\geometry{margin=2.5cm}
\usepackage{graphicx}
\usepackage{hyperref}
\usepackage{listings}
\usepackage{xcolor}
\usepackage{booktabs}
\usepackage{biblatex}
\usepackage{enumitem}

\addbibresource{references.bib}

% Listing styles
\lstdefinestyle{solidity}{
    language=C++,
    basicstyle=\footnotesize\ttfamily,
    keywordstyle=\color{blue},
    commentstyle=\color{gray},
    stringstyle=\color{green!50!black},
    numbers=left,
    numberstyle=\tiny,
    frame=single,
    breaklines=true,
    captionpos=b,
}

\lstdefinestyle{rust}{
    language=Rust,
    basicstyle=\footnotesize\ttfamily,
    keywordstyle=\color{purple},
    commentstyle=\color{gray},
    stringstyle=\color{brown},
    numbers=left,
    numberstyle=\tiny,
    frame=single,
    breaklines=true,
    captionpos=b,
}

\lstdefinestyle{func}{
    language=C,
    basicstyle=\footnotesize\ttfamily,
    keywordstyle=\color{teal},
    commentstyle=\color{gray},
    stringstyle=\color{magenta},
    numbers=left,
    numberstyle=\tiny,
    frame=single,
    breaklines=true,
    captionpos=b,
}

\title{Chain of Contracts: \\ Ethereum, Polkadot, and TON \\\large A Comparative Analysis}
\author{Blockchain Systems Lab --- ChainOfContractsLab}
\date{\today}

\begin{document}

\maketitle
\thispagestyle{empty}

\begin{abstract}
This document provides a comparative analysis of smart contract platforms for
implementing the ``chain of contracts'' pattern --- a design where ContractA
calls ContractB, which calls ContractC in sequence. We analyze three
fundamentally different platforms: Ethereum (synchronous EVM with account
model), Polkadot (synchronous ink! with weight model), and TON (asynchronous
actor model with TVM). For each platform, we examine the execution model,
gas/fee structure, security properties, programming language, and practical
trade-offs. This analysis accompanies the hands-on lab where students implement
the pattern on Ethereum and Polkadot.
\end{abstract}

\tableofcontents
\newpage

\section{Asynchronous vs Synchronous Execution Models}

\subsection{Ethereum: Synchronous EVM Calls}

Ethereum smart contracts execute within a synchronous, single-threaded environment.
When ContractA calls ContractB within the same transaction, the EVM pauses A's
execution, runs B to completion, and returns control to A. This model provides
\textbf{atomicity}: if any step in the chain reverts, the entire transaction
reverts, and all state changes are discarded.

\begin{itemize}
    \item \textbf{Deterministic ordering}: ContractB runs before ContractC,
          and the order is guaranteed by the EVM call stack.
    \item \textbf{Atomic rollback}: A failure in ContractC reverts B and A.
    \item \textbf{Reentrancy risk}: Synchronous calls create reentrancy vectors
          (e.g., the 2016 DAO hack), requiring careful use of checks-effects-interactions.
\end{itemize}

\subsection{Polkadot: Synchronous Cross-Contract Calls (ink! 5.0)}

ink! contracts on Polkadot also execute synchronously within a single dispatch.
When ContractA calls ContractB via \texttt{build\_call}, execution blocks until
B returns. The model closely mirrors Ethereum's in terms of atomicity:

\begin{itemize}
    \item \textbf{Same-block execution}: A $\to$ B $\to$ C all occur in one block.
    \item \textbf{Gas limits per call}: Each cross-contract call can specify its
          own gas limit, preventing one contract from exhausting another's budget.
    \item \textbf{Weight-based pricing}: Polkadot uses a weight system rather
          than gas, but the synchronous semantics are identical for the caller.
\end{itemize}

\subsection{TON: Asynchronous Message Passing}

TON fundamentally differs from Ethereum and Polkadot by adopting an
\textbf{asynchronous actor model}. Smart contracts (actors) communicate
exclusively via messages. When ContractA sends a message to ContractB:

\begin{itemize}
    \item ContractA does \textbf{not} wait for ContractB to process the message.
    \item The message is placed in a queue and may be processed in a different
          shard, by a different validator, in a different block.
    \item ContractA's execution completes immediately; there is no automatic
          rollback if ContractB fails.
    \item Error handling requires explicit \textbf{bounce messages}: if B fails,
          it can send a bounce message back to A, but A must handle it explicitly.
\end{itemize}

\subsection{Impact on the Chain-of-Contracts Pattern}

The choice of execution model directly affects how the A $\to$ B $\to$ C chain
must be designed:

\begin{table}[h]
\centering
\begin{tabular}{|l|c|c|c|}
\hline
\textbf{Property} & \textbf{Ethereum} & \textbf{Polkadot} & \textbf{TON} \\
\hline
Atomicity & Full & Full & None (by default) \\
Error propagation & Automatic revert & Automatic revert & Manual (bounce) \\
Call ordering & Guaranteed & Guaranteed & Not guaranteed \\
Shard-crossing & N/A & N/A & Required \\
\hline
\end{tabular}
\caption{Comparison of execution models for A$\to$B$\to$C chain}
\end{table}

For TON, implementing a reliable chain-of-contracts requires:
\begin{enumerate}
    \item Sending the message from A to B with a bounce flag.
    \item B sending its message to C (with B's remaining balance attached).
    \item C sending a confirmation message back through B to A.
    \item A implementing a timeout mechanism in case messages are lost.
\end{enumerate}

This increased complexity is the cost of TON's scalability advantages:
asynchronous messages enable infinite sharding without cross-shard coordination
overhead at execution time.
\section{Gas and Fee Model Comparison}

\subsection{Ethereum: Sender-Pays Model}

In Ethereum, the transaction \textbf{sender} pays for all computation, storage,
and cross-contract calls in a single upfront gas fee. The entire call chain
(A $\to$ B $\to$ C) consumes gas from the same budget:

\begin{itemize}
    \item Gas is purchased upfront via \texttt{gasLimit} $\times$ \texttt{gasPrice}.
    \item Unused gas is refunded to the sender.
    \item If gas runs out mid-chain, the entire chain reverts (atomicity).
    \item \texttt{EIP-1559} introduced a base fee that is burned and a priority
          fee (tip) for validators.
\end{itemize}

\subsection{Polkadot: Weight-Based Model}

Polkadot does not use gas in the Ethereum sense. Instead, it uses a
\textbf{weight} system:

\begin{itemize}
    \item Each extrinsic declares a \texttt{weight} (computation + storage cost).
    \item Weight is deducted from a per-block weight limit.
    \item Contracts pay for their own execution from their \textbf{own balance},
          not the caller's.
    \item Cross-contract calls can specify separate \texttt{gas\_limit} values
          to prevent one contract from starving another.
    \item Storage is charged via a \textbf{deposit} model: contracts lock a
          refundable deposit for each storage item.
\end{itemize}

For the A $\to$ B $\to$ C chain, each intermediate contract must have sufficient
balance to cover its own execution. ContractA attaches enough gas to cover B's
execution, and B attaches enough to cover C's.

\subsection{TON: Separate Balances and Rent}

TON's fee model introduces several unique features that significantly impact
chain-of-contracts design:

\begin{itemize}
    \item \textbf{Each contract has its own balance.} There is no "sender pays
          all" model. ContractA must attach enough TON to its outbound message
          to cover ContractB's execution. This is called the \texttt{gram}
          attached to the message.
    \item \textbf{Forward fees}: Each hop in the chain incurs a forwarding fee
          paid by the sender of that message. ContractA pays to reach B, B pays
          to reach C.
    \item \textbf{Storage rent}: TON contracts pay rent for their persistent
          state. A contract with insufficient balance may be frozen or deleted.
          Long-running chains must account for ongoing storage costs.
    \item \textbf{Action fees}: Each outbound action (sending a message) costs
          fees based on the message size and destination shard.
    \item \textbf{No automatic refunds}: Unlike Ethereum, there is no gas refund
          mechanism. Unused attached value stays with the recipient contract.
\end{itemize}

\subsection{Practical Implications for the Chain Pattern}

\begin{table}[h]
\centering
\begin{tabular}{|l|c|c|c|}
\hline
\textbf{Aspect} & \textbf{Ethereum} & \textbf{Polkadot} & \textbf{TON} \\
\hline
Who pays? & Sender & Each contract & Message sender at each hop \\
Chain failure cost & Full (revert) & Partial (weight) & Accumulated (fees spent) \\
Storage cost & Per byte (one-time) & Deposit (refundable) & Rent (ongoing) \\
Refund mechanism & Yes & No & No \\
\hline
\end{tabular}
\caption{Gas/fee model comparison for A$\to$B$\to$C}
\end{table}

The key takeaway for students: in TON, each hop in the chain is an independent
economic transaction. Attaching insufficient value to a message will cause the
chain to break mid-way, and unlike Ethereum, there is no automatic rollback.
This makes TON's model more complex but also more flexible -- intermediate
contracts are not dependent on the sender's generosity for their execution
budget.
\section{Actor Model vs Account Model}

\subsection{Ethereum and Polkadot: Shared Global State}

Both Ethereum and Polkadot use an \textbf{account model} where smart contracts
operate on shared, globally-accessible state. In Ethereum, any contract can read
the storage of any other contract (via \texttt{SLOAD}). In Polkadot, ink!
contracts are more isolated by default, but the underlying execution environment
is still a shared state machine.

Key characteristics:
\begin{itemize}
    \item \textbf{Reentrancy is possible}: Because the call stack is synchronous,
          a contract can be called again while it is still executing (reentrancy).
    \item \textbf{Global composability}: ContractA can directly call any function
          on ContractB without going through a message queue.
    \item \textbf{MEV (Maximal Extractable Value)}: The shared mempool allows
          validators to reorder transactions, potentially extracting value from
          the chain.
    \item \textbf{Deterministic execution}: All nodes replay the same sequence
          of operations and converge on the same state.
\end{itemize}

\subsection{TON: The Actor Model}

TON implements the \textbf{Actor model}, originally from distributed computing
(Hewitt, Bishop, and Steiger, 1973). Each smart contract is an \textbf{actor}
with:

\begin{itemize}
    \item \textbf{Exclusive private state}: No other contract can read or modify
          a contract's storage directly. State is only accessible through the
          contract's own message handlers.
    \item \textbf{Message-based communication}: Actors communicate exclusively
          via asynchronous messages. There are no direct function calls across
          contract boundaries.
    \item \textbf{No reentrancy}: Since there is no synchronous call stack,
          reentrancy is structurally impossible. A contract cannot be interrupted
          mid-execution by another call.
    \item \textbf{Each actor processes one message at a time}: The TON runtime
          ensures serial message processing per contract, eliminating race
          conditions within a single actor.
\end{itemize}

\subsection{Security Implications for the Chain Pattern}

\begin{table}[h]
\centering
\begin{tabular}{|l|c|c|c|}
\hline
\textbf{Security Property} & \textbf{Ethereum} & \textbf{Polkadot} & \textbf{TON} \\
\hline
Reentrancy risk & High & Medium & None \\
Cross-contract read & Yes & Limited & No \\
Atomic execution & Yes & Yes & No \\
Message ordering & N/A & N/A & Non-deterministic \\
Parallelism & None & None & Shard-level \\
\hline
\end{tabular}
\caption{Security model comparison}
\end{table}

For our A $\to$ B $\to$ C chain pattern:

\begin{itemize}
    \item \textbf{Ethereum}: The chain is atomic and synchronous, but requires
          reentrancy guards (\texttt{ReentrancyGuard} or checks-effects-interactions).
          A malicious ContractB could attempt to re-enter ContractA during the call.
    \item \textbf{Polkadot}: Similar security profile to Ethereum, but ink!'s
          type system and the \texttt{build\_call} API make accidental reentrancy
          less likely.
    \item \textbf{TON}: Reentrancy is impossible, but the chain is non-atomic.
          The professor's grader must check for bounce messages or confirmation
          receipts. If ContractC receives the message but ContractB crashes before
          returning a confirmation, ContractA never knows the chain completed.
\end{itemize}

\subsection{State Isolation and Composability Trade-off}

TON's actor model provides stronger security guarantees (no reentrancy, no
unauthorized state reads) at the cost of composability. In Ethereum, composing
contracts (using the output of one contract as input to another in the same
transaction) is seamless. In TON, such composition requires multi-step
asynchronous protocols, which are more complex to design and verify.

For educational purposes, this trade-off is instructive: students should
understand that stronger isolation means weaker composition, and that the choice
of platform determines which properties are guaranteed by the runtime versus
enforced by the developer.
\section{Language Comparison: Solidity vs ink! vs FunC}

\subsection{Ethereum: Solidity}

Solidity is the dominant smart contract language for EVM-compatible chains. It
is a high-level, curly-bracket language influenced by C++, Python, and
JavaScript.

\begin{lstlisting}[language=Solidity, caption=Chain forwarder in Solidity, label=lst:solidity]
function validateAndForward(
    uint256 chainId,
    bytes calldata data
) external onlyContractA returns (bool) {
    bytes memory transformed = abi.encodePacked(
        "[VALIDATED]", data
    );
    emit MessageValidated(chainId, true, transformed);
    return IContractC(contractC).receiveFinal(
        chainId, transformed
    );
}
\end{lstlisting}

Key traits:
\begin{itemize}
    \item \textbf{Familiar syntax}: Low barrier for students with C/Java background.
    \item \textbf{Mature ecosystem}: Hardhat, Foundry, Remix -- extensive tooling.
    \item \textbf{Interface-based composition}: Cross-contract calls use interface types.
    \item \textbf{Gas-aware design}: Students must optimize storage and computation.
\end{itemize}

\subsection{Polkadot: ink! 5.0}

ink! is a Rust-based eDSL (embedded Domain-Specific Language) for WebAssembly
smart contracts on Substrate-based chains.

\begin{lstlisting}[language=Rust, caption=Chain forwarder in ink! 5.0, label=lst:ink]
fn validate_and_forward(
    &mut self,
    chain_id: u32,
    data: Vec<u8>,
) -> bool {
    let mut transformed = b"[VALIDATED]".to_vec();
    transformed.extend_from_slice(&data);

    self.env().emit_event(MessageValidated {
        chain_id, valid: true,
        transformed_data: transformed.clone(),
    });

    build_call::<DefaultEnvironment>()
        .call(self.contract_c.unwrap())
        .exec_input(
            ExecutionInput::new(
                Selector::new(
                    ink::selector_bytes!("receive_final")
                )
            )
            .push_arg(chain_id)
            .push_arg(transformed),
        )
        .returns::<bool>()
        .invoke()
}
\end{lstlisting}

Key traits:
\begin{itemize}
    \item \textbf{Rust safety}: Ownership, borrowing, and type safety at compile time.
    \item \textbf{Macro-based}: \texttt{\#[ink(storage)]}, \texttt{\#[ink(message)]} annotations.
    \item \textbf{No standard library}: \texttt{no\_std} environment -- no heap by default.
    \item \textbf{Steeper learning curve}: Students need Rust familiarity.
\end{itemize}

\subsection{TON: FunC}

FunC is a low-level language for TON smart contracts. It compiles to Fift
assembler and then to TVM bytecode.

\begin{lstlisting}[language=C, caption=Simple forwarder in FunC, label=lst:func]
() recv_internal(int msg_value, cell in_msg_cell, slice in_msg_body) impure {
    ;; Parse incoming message
    slice ds = in_msg_body;
    int op = ds~load_uint(32);
    int chain_id = ds~load_uint(32);
    slice data = ds~load_dict_slice();

    ;; Validate (non-empty check)
    throw_if(1001, slice_empty?(data));

    ;; Build forward message to ContractB
    cell forward_msg = begin_cell()
        .store_uint(0xba5eba11, 32)
        .store_uint(chain_id, 32)
        .store_slice(data)
        .end_cell();

    ;; Send with bounce flag and remaining balance
    send_raw_message(forward_msg, 64);

    ;; Emit event via storage
    cell event = begin_cell()
        .store_uint(0xcafe, 16)
        .store_uint(chain_id, 32)
        .end_cell();
    set_data(begin_cell().store_ref(event).end_cell());
}
\end{lstlisting}

Key traits:
\begin{itemize}
    \item \textbf{Stack-based}: FunC operates on a stack machine (TVM).
    \item \textbf{Assembly-like}: Manual cell manipulation and gas management.
    \item \textbf{Low-level control}: Full control over message layout and fees.
    \item \textbf{Steepest learning curve}: Most challenging for students.
\end{itemize}

\subsection{Comparative Analysis for Education}

\begin{table}[h]
\centering
\begin{tabular}{|l|c|c|c|}
\hline
\textbf{Criterion} & \textbf{Solidity} & \textbf{ink! 5.0} & \textbf{FunC} \\
\hline
Learning curve & Low & Medium & High \\
Type safety & Low (runtime) & High (compile-time) & Low \\
Tooling maturity & Excellent & Good & Emerging \\
Cross-contract calls & Trivial & Moderate & Complex \\
Debugging & Hardhat/Foundry & Contracts UI & TON CLI \\
Gas estimation & Built-in & Weight system & Manual \\
\hline
\end{tabular}
\caption{Language comparison for smart contract education}
\end{table}

For a blockchain course, the recommended progression is:
\begin{enumerate}
    \item Start with Solidity (familiar syntax, abundant resources).
    \item Move to ink! (safety lessons from Rust, different cross-contract API).
    \item End with FunC (understanding low-level execution, message cells,
          and the actor model at the VM level).
\end{enumerate}

This progression takes students from high-level abstraction to low-level control,
building a complete mental model of how different blockchain VMs approach the
same problem.
\section{Practical Trade-offs and Recommendations}

\subsection{Decision Matrix}

The following table summarizes the key differences across all three platforms
for the chain-of-contracts pattern:

\begin{table}[h]
\centering
\small
\begin{tabular}{|l|c|c|c|}
\hline
\textbf{Factor} & \textbf{Ethereum} & \textbf{Polkadot} & \textbf{TON} \\
\hline
Atomic chain execution & Yes & Yes & No \\
Reentrancy protection & Manual & Built-in (mostly) & Structural \\
Cross-contract complexity & Low & Medium & High \\
Event verification & Trivial & Moderate & Indirect \\
Max sharding & None & Para-chain & Infinite \\
User base & Largest & Growing & Niche \\
Tooling maturity & Mature & Maturing & Early \\
Gas model clarity & Simple & Moderate & Complex \\
Storage cost model & One-time & Deposit & Rent \\
\hline
\end{tabular}
\caption{Comprehensive comparison across platforms}
\end{table}

\subsection{When to Use Each Platform}

\begin{itemize}
    \item \textbf{Ethereum} is the best choice for:
    \begin{itemize}
        \item Rapid prototyping and student education.
        \item Atomic chains where intermediate state must not be persisted on failure.
        \item Applications requiring maximum developer tooling and community support.
        \item DeFi protocols where synchronous composability is essential.
    \end{itemize}

    \item \textbf{Polkadot (ink!)} is the best choice for:
    \begin{itemize}
        \item Projects requiring cross-chain interoperability via XCM.
        \item Applications needing formal verification (Rust's type system helps).
        \item Teams already familiar with Rust.
        \item Scenarios where upgradeable smart contracts (via \texttt{set\_code\_hash}) are needed.
    \end{itemize}

    \item \textbf{TON} is the best choice for:
    \begin{itemize}
        \item High-throughput applications serving millions of users.
        \item Systems where sharding is required today, not as a future roadmap item.
        \item Applications that can tolerate asynchronous message passing.
        \item Projects that benefit from TON's native Telegram integration.
    \end{itemize}
\end{itemize}

\subsection{Hybrid Approaches}

For production systems, platforms can be combined:

\begin{itemize}
    \item \textbf{ETH + TON bridges}: Use Ethereum for the atomic chain core
          and TON for high-volume message relaying to end users.
    \item \textbf{Polkadot + TON}: Deploy the chain-of-contracts as a parachain
          (synchronous, sovereign) and bridge messages to TON for distribution.
    \item \textbf{TON with bounce guards}: Implement the A $\to$ B $\to$ C chain
          in TON using bounce messages and confirmation timeouts, accepting
          eventual consistency for scalability.
\end{itemize}

\subsection{Educational Recommendations}

For a chain-of-contracts lab, we recommend:

\begin{enumerate}
    \item \textbf{Implement on Ethereum first} -- the synchronous model is
          easier to understand and debug. Students can verify the chain with
          a single transaction receipt.
    \item \textbf{Port to Polkadot} -- demonstrates how the same pattern
          translates to a different VM with stronger type safety and a
          different cross-contract API.
    \item \textbf{Analyze TON theoretically (this document)} -- understand
          how the pattern must be redesigned for the actor model, without
          requiring FunC proficiency.
\end{enumerate}

This three-step approach gives students both practical implementation skills
and a theoretical framework for evaluating smart contract platforms.

\newpage
\printbibliography

\end{document}